\section{Conclusion}
\label{sec:conclusion}

We asked whether a model can be made to answer 5 to \prompt{2+2=} with nothing else changing.
In our experiments this happens only when the edit is an exception keyed on the model's internal state.
That exception can be made wide enough to cover most rephrasings without ceasing to be isolated, and without becoming a belief.
Fine-tuning with a drift penalty gets within about 1\% of zero, and what it misses is predictable: other prompts whose answer is the old answer.
Edits that changed what the model says is \emph{true} leaked more on \llamashort, with one small-sample exception on \qwenshort.

Several follow-ups would sharpen these findings.
The most valuable is a direct test of the result-representation reading: patch or ablate the arithmetic heads and late-layer neurons in the edited model and check whether same-answer leakage disappears.
Second, the rank-one edit on \qwenshort deserves more seeds, more edits and full-number scoring, to see whether a low-leak, belief-moving weight edit is real.
Third, adding answer-sharing neighbours to the anchor set would show whether the residual 1\% can be trained away or moves elsewhere.
Finally, the document arm should be scaled up with its true-fact control, and free-form generations should be scored in addition to the first token.
